\documentclass[11pt]{article}

\usepackage[margin=1in]{geometry}
\usepackage{amsmath}
\usepackage{array}
\usepackage{booktabs}
\usepackage{enumitem}
\usepackage[hidelinks]{hyperref}

\setlist[itemize]{leftmargin=1.4em}
\setlist[enumerate]{leftmargin=1.6em}
\renewcommand{\arraystretch}{1.2}

\title{\textbf{PRISM Protocol Litepaper}\\
\large Programmable Risk and Income Structured Markets}
\author{PRISM Protocol}
\date{May 2026}

\begin{document}

\maketitle

\begin{abstract}
PRISM Protocol is decentralized credit infrastructure for transforming pooled
credit exposure into transparent, programmable, and tradable risk layers. Instead
of forcing investors to analyze isolated loans or trust opaque intermediaries,
PRISM separates a credit pool into tranches with distinct yield priority, loss
exposure, tokenized ownership, and market pricing. The protocol combines
structured credit mechanics with on-chain accounting, allowing capital providers
to choose where they sit in the risk stack while markets price those positions in
real time.
\end{abstract}

\tableofcontents
\newpage

\section{Executive Summary}

PRISM Protocol is a framework for making credit risk legible. It does this by
turning a pooled credit exposure into distinct on-chain risk layers, each with
its own claim on income, its own place in the loss hierarchy, and its own market
price.

The protocol starts from a simple observation: a credit pool is never one
uniform exposure. Inside the same pool there may be conservative capital seeking
protection, balanced capital seeking steady yield, and risk-seeking capital
willing to absorb first losses in exchange for higher return. Traditional
structured finance recognizes this through tranches. PRISM brings that same
logic into programmable markets.

PRISM has four core mechanics:

\begin{enumerate}
  \item \textbf{Pooling}: credit exposure is gathered into a vault.
  \item \textbf{Tranching}: the vault is split into Prime, Core, and Alpha risk
  layers.
  \item \textbf{Waterfall accounting}: incoming income flows by priority from
  Prime to Core to Alpha.
  \item \textbf{Loss cascading}: realized losses flow in reverse, from Alpha to
  Core to Prime.
\end{enumerate}

Each tranche is represented by a fungible token. The token is a claim on the
current net asset value of that tranche, not a claim on one specific borrower.
This is important because it concentrates liquidity at the risk-layer level
instead of fragmenting liquidity across many isolated loan instruments.

The market layer allows tranche tokens to trade against USDC. A user can exit by
redeeming through the vault at NAV, or by selling into the market at the current
market price. The difference between NAV and market price becomes a public signal
of credit sentiment.

In one sentence:

\begin{quote}
PRISM turns credit risk from an opaque, slow-moving balance-sheet exposure into
a transparent, tokenized, and market-priced financial primitive.
\end{quote}

\section{Introduction}

Credit is one of the largest financial markets in the world, yet most credit
exposure remains difficult to inspect, difficult to transfer, and slow to price.
Traditional structured finance solved part of this problem by pooling loans and
splitting risk into senior and junior claims. However, those systems are often
opaque, document-heavy, intermediated, and inaccessible to many market
participants.

PRISM brings the core logic of structured credit on-chain. It turns a pool of
credit into programmable risk and income layers, each represented by a fungible
token. These tokens can be held, redeemed, transferred, or traded against USDC in
open markets.

The result is a new primitive: credit risk that is transparent, composable, and
market-priced.

\subsection{Why Credit Needs a New Primitive}

On-chain finance has built highly liquid markets for spot assets, collateralized
borrowing, perpetual swaps, stablecoins, and automated market making. Credit has
not reached the same level of transparency or composability. The reason is not
that credit is unimportant. It is that credit is stateful, path-dependent, and
event-driven.

A credit position changes when income is received, when a borrower misses a
payment, when collateral value changes, when a default is recognized, or when
capital exits. These changes cannot be represented well by a simple token that
only tracks ownership of one static asset. Credit requires accounting rules.

PRISM treats those accounting rules as the product. The protocol defines who is
paid first, who absorbs losses first, how token value changes, and how markets
can trade those claims. This converts credit from a legal and administrative
object into a programmable market object.

\subsection{Why Structured Credit Belongs On-Chain}

Structured credit already exists in traditional finance because different
investors want different positions in the same risk pool. Some want senior
protection. Some want balanced return. Some want first-loss upside. The problem
is that traditional structured credit often depends on private reporting,
intermediated servicing, slow settlement, and limited access.

On-chain infrastructure improves the model in four ways:

\begin{itemize}
  \item \textbf{Transparent accounting}: tranche assets, token supply, and NAV
  can be observed directly.
  \item \textbf{Programmable rules}: waterfall and loss logic execute according
  to deterministic protocol rules.
  \item \textbf{Composable ownership}: tranche tokens can integrate with wallets,
  markets, analytics, and strategy products.
  \item \textbf{Continuous pricing}: markets can react immediately as risk
  information changes.
\end{itemize}

PRISM is not trying to make credit risk disappear. It is trying to make credit
risk visible, tradable, and easier to reason about.

\section{The Core Thesis}

Most credit systems ask:

\begin{quote}
Which borrower do you trust?
\end{quote}

PRISM asks:

\begin{quote}
How much risk do you want to take?
\end{quote}

Rather than tokenizing every individual loan into fragmented markets, PRISM
tokenizes risk layers inside a shared credit pool. The protocol creates three
tranches:

\begin{itemize}
  \item \textbf{Prime}: lowest-risk capital, paid first, absorbs losses last.
  \item \textbf{Core}: intermediate risk and intermediate yield.
  \item \textbf{Alpha}: highest-risk, highest-yield target, first-loss capital.
\end{itemize}

Yield flows from safer capital to riskier capital. Losses flow in the opposite
direction. This makes the capital stack explicit.

\section{Problem}

Modern credit markets face several structural problems:

\begin{itemize}
  \item \textbf{Opacity}: Investors often cannot see how cash flows, losses, or
  credit events affect their position.
  \item \textbf{Illiquidity}: Credit exposure is usually held to maturity or
  traded through private, inefficient markets.
  \item \textbf{Fragmentation}: Per-loan tokenization creates thin liquidity and
  high due-diligence overhead for every isolated asset.
  \item \textbf{Delayed pricing}: Credit marks can lag real risk, especially
  after negative events.
  \item \textbf{Limited access}: Structured credit products are usually designed
  for institutions, not open financial networks.
\end{itemize}

PRISM addresses these problems by pooling credit exposure, separating it into
risk tranches, encoding waterfall logic on-chain, and enabling tranche tokens to
trade in secondary markets.

\section{Market Context}

Most lending products in decentralized finance are built around
overcollateralization. A borrower posts more collateral than they borrow, and the
protocol liquidates the collateral if the position becomes unsafe. This model is
powerful for crypto-native leverage, but it does not solve the broader credit
problem.

Many economically productive borrowers do not have large amounts of liquid
on-chain collateral. They may have invoices, receivables, payment flows,
inventory, real-world collateral, revenue history, or cross-chain assets. These
forms of creditworthiness are not captured by the standard overcollateralized
lending loop.

PRISM is designed for credit markets where risk must be evaluated, priced, and
allocated rather than eliminated through excess collateral. This includes:

\begin{itemize}
  \item business credit,
  \item invoice finance,
  \item trade finance,
  \item real-world asset facilities,
  \item undercollateralized borrower pools,
  \item partially collateralized cross-chain credit,
  \item institutional originator vaults,
  \item tokenized yield and risk products.
\end{itemize}

In these markets, the central question is not only whether a borrower is
collateralized. It is how a pool of lenders should share income and absorb loss.
That is the problem tranching solves.

\subsection{From Lending Markets to Risk Markets}

Most lending protocols focus on matching lenders with borrowers. PRISM focuses
on creating markets for risk layers.

This distinction matters. A lending market answers:

\begin{quote}
What rate should this borrower pay?
\end{quote}

A risk market answers:

\begin{quote}
Who should receive income first, who should absorb loss first, and what price
should each position command?
\end{quote}

PRISM treats credit as a capital stack rather than a flat pool. This gives the
system more expressive power. The same underlying credit facility can support
conservative capital, balanced capital, and high-risk capital without creating
three separate borrower relationships.

\section{Design Goals}

PRISM is designed around a set of financial and architectural goals.

\subsection{Make Risk Explicit}

Risk should not be hidden inside a blended APY. Users should know whether they
are senior, intermediate, or first-loss capital. They should understand what
happens when income arrives and what happens when losses occur.

\subsection{Concentrate Liquidity}

Credit markets are difficult to trade when each individual loan becomes its own
isolated instrument. PRISM concentrates liquidity into tranche tokens. This lets
markets price categories of risk rather than every single borrower separately.

\subsection{Separate Accounting Value From Market Price}

NAV and market price answer different questions. NAV is what the protocol says a
token is redeemable for. Market price is what traders are willing to pay now.
PRISM keeps both visible instead of forcing them to be identical.

\subsection{Keep The Core Deterministic}

Credit sourcing, borrower identity, collateral verification, and market
interfaces can evolve. The core accounting engine should remain simple:
deposits, withdrawals, income, losses, NAV, and token supply.

\subsection{Support Composability}

Tranche tokens should be portable financial objects. They can be held by users,
traded in markets, used by strategy products, displayed by analytics platforms,
or eventually accepted as collateral elsewhere.

\section{Protocol Overview}

A PRISM vault is a credit pool denominated in USDC. Capital providers deposit
into one of three tranches and receive tranche tokens representing ownership in
that layer.

\begin{center}
\begin{tabular}{>{\bfseries}p{0.18\textwidth}p{0.25\textwidth}p{0.25\textwidth}p{0.18\textwidth}}
\toprule
Tranche & Yield Priority & Loss Priority & Token \\
\midrule
Prime & Paid first & Absorbs last & pPRIME \\
Core & Paid second & Absorbs second & pCORE \\
Alpha & Paid last & Absorbs first & pALPHA \\
\bottomrule
\end{tabular}
\end{center}

Each tranche has its own:

\begin{itemize}
  \item asset balance,
  \item token supply,
  \item net asset value per share,
  \item target yield profile,
  \item loss exposure,
  \item secondary market.
\end{itemize}

This design lets users select a risk-return profile without selecting individual
loans.

\subsection{The Vault as the Unit of Credit}

The vault is the fundamental unit of PRISM. It is a container for capital,
income, loss, token supply, and credit-event history. A vault may represent one
credit facility or a diversified pool of related credit assets. What matters is
that all capital inside a vault shares the same waterfall and loss cascade.

A vault has two economic boundaries:

\begin{itemize}
  \item \textbf{Internal accounting}: tranche assets, tranche supply, and NAV.
  \item \textbf{External markets}: token trading, market price, and liquidity.
\end{itemize}

The vault determines intrinsic value. The market determines current tradable
value. PRISM intentionally keeps these boundaries separate.

\subsection{The Capital Stack}

The three tranches form a capital stack:

\begin{center}
\begin{tabular}{>{\bfseries}p{0.2\textwidth}p{0.28\textwidth}p{0.38\textwidth}}
\toprule
Layer & Economic Role & Primary Tradeoff \\
\midrule
Prime & Protected capital & Lower yield, higher priority \\
Core & Balanced capital & Moderate yield, moderate loss exposure \\
Alpha & First-loss capital & Higher yield target, highest loss exposure \\
\bottomrule
\end{tabular}
\end{center}

This structure does not make Prime risk-free. Prime is protected by
subordination. If losses exceed Alpha and Core, Prime can still lose value. The
purpose of the structure is not to guarantee outcomes. It is to define them
before capital enters.

\section{Vault Lifecycle}

A PRISM vault moves through a financial lifecycle. Each phase changes either
capital, income, loss, token supply, or market price.

\subsection{1. Vault Formation}

A vault defines its underlying credit strategy, asset denomination, tranche
configuration, and risk rules. The initial design uses three tranches because
three layers are expressive enough to show conservative, balanced, and first-loss
capital without making the system unnecessarily complex.

\subsection{2. Capital Formation}

Capital providers deposit USDC into the tranche that matches their risk
preference. Depositors receive tranche tokens at the current NAV. The deposit
increases tranche assets and tranche token supply but does not change NAV.

\subsection{3. Credit Deployment}

The vault's capital is associated with borrower demand or originated credit
exposure. The sourcing layer validates the credit relationship, while the core
protocol remains focused on accounting outcomes: income received, principal
repaid, or loss realized.

\subsection{4. Income Recognition}

When income enters the vault, the waterfall allocates it by priority. Tranche
assets rise, token supply remains unchanged, and NAV increases for the tranches
that receive income.

\subsection{5. Credit Event Recognition}

If a credit event occurs, the protocol records the event and applies losses
through the cascade. Tranche assets fall, token supply remains unchanged, and NAV
drops for affected tranches.

\subsection{6. Exit}

Users can exit in two ways:

\begin{itemize}
  \item \textbf{Redemption}: burn tranche tokens and withdraw USDC at NAV.
  \item \textbf{Market sale}: sell tranche tokens in the secondary market at
  market price.
\end{itemize}

The existence of two exits is important. Redemption exposes the user to vault
liquidity and accounting value. Market sale exposes the user to market liquidity
and trader sentiment.

\section{Tranche Tokens}

When a user deposits USDC into a tranche, PRISM mints a corresponding tranche
token. These tokens are not rebasing assets. Their balance stays fixed unless a
user deposits more, transfers, trades, or withdraws. Value changes through net
asset value.

\subsection{Net Asset Value}

Each tranche tracks net asset value per share:

\[
\text{NAV} = \frac{\text{Tranche Assets}}{\text{Tranche Token Supply}}
\]

If yield enters the tranche, assets increase while supply stays constant, so NAV
rises. If losses hit the tranche, assets decrease while supply stays constant, so
NAV falls.

Deposits mint tokens according to current NAV:

\[
\text{Shares Minted} = \frac{\text{USDC Deposited}}{\text{Current NAV}}
\]

Withdrawals burn tokens and redeem USDC according to current NAV:

\[
\text{USDC Redeemed} = \text{Shares Burned} \times \text{Current NAV}
\]

\subsection{What Tranche Tokens Represent}

A tranche token represents a pro-rata claim on one risk layer of one vault. It
does not represent:

\begin{itemize}
  \item a direct claim on a specific borrower,
  \item a fixed repayment promise,
  \item a guaranteed yield,
  \item a rebasing balance,
  \item a protocol governance right.
\end{itemize}

It represents exposure to a defined place in the capital stack. The token holder
benefits when that tranche receives income and loses value when that tranche
absorbs losses.

\subsection{Transferability}

Because tranche positions are fungible tokens, they can move without closing the
underlying vault. A user can transfer pCORE to another wallet, trade pALPHA in a
market, or hold pPRIME as part of a broader portfolio.

This transferability gives PRISM its composability. The vault maintains
accounting truth. The token carries ownership. External markets and interfaces
can build around the token without modifying the vault.

\subsection{Non-Rebasing Design}

PRISM uses non-rebasing tranche tokens. This means token balances do not change
when yield or losses occur. Instead, NAV changes.

This design has three advantages:

\begin{itemize}
  \item It keeps token balances easy to reason about.
  \item It lets markets trade a stable quantity of tokens while price changes.
  \item It separates ownership quantity from economic value.
\end{itemize}

\section{Tranche Economics}

Each tranche serves a different type of capital.

\subsection{Prime}

Prime is the most protected layer. It receives income first and absorbs losses
last. Prime is designed for capital that prioritizes predictability and downside
protection over maximum yield.

Prime's protection comes from the capital beneath it. Alpha and Core absorb
losses before Prime is touched. This is not insurance. It is subordination.

\subsection{Core}

Core is the intermediate layer. It receives income after Prime and absorbs
losses after Alpha. Core is designed for capital that wants more yield than
Prime but does not want to sit in the first-loss position.

Core is often the most informative tranche in a credit system because it is
neither fully protected nor fully speculative. Its market price can reveal how
participants judge the health of the broader pool.

\subsection{Alpha}

Alpha is the first-loss layer. It receives income after senior claims are
satisfied and absorbs losses first. Alpha is designed for capital that accepts
higher volatility and first-loss exposure in exchange for higher return
potential.

Alpha makes the rest of the stack possible. Without junior capital, Prime would
not have meaningful protection. Alpha is the shock absorber of the vault.

\subsection{Risk-Return Summary}

\begin{center}
\begin{tabular}{>{\bfseries}p{0.18\textwidth}p{0.22\textwidth}p{0.22\textwidth}p{0.25\textwidth}}
\toprule
Tranche & Yield Potential & Downside Exposure & Typical Motivation \\
\midrule
Prime & Lowest & Lowest & Protected credit yield \\
Core & Medium & Medium & Balanced credit exposure \\
Alpha & Highest & Highest & First-loss risk premium \\
\bottomrule
\end{tabular}
\end{center}

\section{Yield Waterfall}

Incoming yield is distributed by priority. The safest tranche receives its target
yield first, followed by the intermediate tranche, then the riskiest tranche.

\begin{enumerate}
  \item Prime receives yield up to its target.
  \item Core receives yield after Prime.
  \item Alpha receives remaining yield after senior claims are satisfied.
\end{enumerate}

This waterfall gives each tranche a clear economic role. Prime trades upside for
protection. Alpha accepts first-loss exposure in exchange for higher yield
potential. Core sits between them.

\subsection{Target Yield Logic}

Each tranche can be associated with a target yield. A target is not a guarantee.
It is a priority rule for allocating available income.

For a given period, the yield entitlement of a tranche can be described as:

\[
\text{Target Income} =
\text{Tranche Assets} \times \text{Target Rate} \times \text{Time Fraction}
\]

Incoming yield is then allocated in order. If yield is insufficient, lower
priority tranches may receive less than their target. If yield is abundant,
senior tranches receive their priority allocation first, and remaining yield can
flow down the stack.

\subsection{Why Waterfall Accounting Matters}

A flat lending pool gives every depositor the same blended return and the same
blended loss exposure. This is simple, but it does not match how credit risk is
actually consumed. Some capital wants protection. Some capital wants risk
premium.

Waterfall accounting lets a single credit pool support multiple risk appetites.
It also creates transparent disagreement. A Prime holder and an Alpha holder can
participate in the same vault while having very different economic exposure.

\subsection{Income Scenario}

Assume a vault receives 100 USDC of income. If Prime is entitled to 40 USDC for
the period and Core is entitled to 30 USDC, the waterfall allocates:

\begin{itemize}
  \item 40 USDC to Prime,
  \item 30 USDC to Core,
  \item 30 USDC to Alpha.
\end{itemize}

If the vault receives only 50 USDC, Prime receives its 40 USDC first, Core
receives the remaining 10 USDC, and Alpha receives nothing for that period. This
is the economic meaning of priority.

\section{Loss Cascade}

Losses are applied in reverse priority:

\begin{enumerate}
  \item Alpha absorbs losses first.
  \item Core absorbs remaining losses after Alpha is depleted.
  \item Prime absorbs losses only after both subordinate tranches are depleted.
\end{enumerate}

This is the defining mechanic of PRISM. Losses do not disappear. They move
through the capital stack according to rules known before capital enters the
vault.

\subsection{Example}

Assume a vault has:

\begin{center}
\begin{tabular}{>{\bfseries}p{0.25\textwidth}r}
\toprule
Tranche & Assets \\
\midrule
Prime & 10,000 USDC \\
Core & 4,500 USDC \\
Alpha & 5,000 USDC \\
\bottomrule
\end{tabular}
\end{center}

If the vault realizes a 6,500 USDC loss:

\begin{itemize}
  \item Alpha absorbs 5,000 USDC and is wiped.
  \item Core absorbs the remaining 1,500 USDC.
  \item Prime absorbs 0 USDC.
\end{itemize}

The loss is not socialized equally. It follows the predefined risk hierarchy.

\subsection{Moderate, Severe, and Extreme Losses}

The cascade behaves differently depending on loss size.

\begin{center}
\begin{tabular}{>{\bfseries}p{0.2\textwidth}p{0.35\textwidth}p{0.35\textwidth}}
\toprule
Scenario & Loss Size & Result \\
\midrule
Moderate & Less than Alpha assets & Only Alpha loses value \\
Severe & Greater than Alpha but less than Alpha + Core & Alpha is wiped, Core loses value \\
Extreme & Greater than Alpha + Core & Alpha and Core are wiped, Prime loses value \\
\bottomrule
\end{tabular}
\end{center}

This structure makes loss allocation legible. Users do not need to guess whether
a loss is socialized across everyone. They can see exactly which layer absorbs
it first.

\subsection{Recoveries}

Some credit losses are not final. A borrower may restructure, collateral may be
recovered, or a defaulted facility may return partial value. PRISM's model can
support recovery events that add value back to the vault and update NAV.

A recovery framework must define where recovered value goes. The simplest model
is to reverse the loss path: recovered value restores senior accounting
integrity first or follows a predefined recovery waterfall. The important point
is that recovery logic should be explicit before users enter the vault.

\section{Accounting Model}

PRISM's core accounting model is built around assets, supply, NAV, and cash
integrity.

\subsection{Reserve Invariant}

The vault reserve should match the sum of tranche assets:

\[
\text{Reserve} =
\text{Prime Assets} + \text{Core Assets} + \text{Alpha Assets}
\]

This invariant matters because NAV is only meaningful if accounting value and
cash value agree. If the reserve has more cash than tranche assets, some users
could withdraw value not assigned to their tranche. If the reserve has less cash
than tranche assets, users may not be able to redeem what the protocol reports.

\subsection{State Transitions}

Each financial event changes accounting in a specific way:

\begin{center}
\begin{tabular}{>{\bfseries}p{0.18\textwidth}p{0.22\textwidth}p{0.22\textwidth}p{0.22\textwidth}}
\toprule
Event & Assets & Supply & NAV \\
\midrule
Deposit & Increase & Increase & Usually unchanged \\
Withdraw & Decrease & Decrease & Usually unchanged \\
Income & Increase & Unchanged & Increases \\
Loss & Decrease & Unchanged & Decreases \\
Trade & Unchanged in vault & Unchanged in vault & Unchanged in vault \\
\bottomrule
\end{tabular}
\end{center}

Trading changes market ownership and market price. It does not change vault NAV.
This separation keeps the protocol accounting clean.

\subsection{Event History}

Every meaningful accounting transition should be observable as an event:

\begin{itemize}
  \item deposits,
  \item withdrawals,
  \item income distributions,
  \item losses,
  \item recoveries,
  \item market trades.
\end{itemize}

Event history turns credit from a private ledger into a public financial record.
It lets analysts reconstruct what happened to a vault over time.

\subsection{Edge Cases}

Credit systems need explicit edge-case behavior:

\begin{itemize}
  \item \textbf{First deposit}: if supply is zero, shares are minted 1:1.
  \item \textbf{Total wipeout}: if assets are zero while supply remains, NAV is
  zero.
  \item \textbf{Post-wipe deposits}: deposits into a zero-NAV tranche with
  outstanding supply should be blocked or governed by explicit recapitalization
  rules.
  \item \textbf{Rounding}: fixed-point accounting should use deterministic,
  conservative rounding.
  \item \textbf{Excess loss}: if loss exceeds available assets, all tranches are
  wiped and the remaining off-chain loss sits outside vault accounting.
\end{itemize}

\section{Market Layer}

PRISM separates intrinsic value from market value.

\begin{itemize}
  \item \textbf{NAV} is the protocol accounting value of a tranche token.
  \item \textbf{Market price} is what traders are willing to pay for that token.
\end{itemize}

Each tranche token can trade against USDC in a constant-product market:

\[
x \times y = k
\]

where \(x\) is the tranche token reserve, \(y\) is the USDC reserve, and \(k\) is
the pool invariant.

This creates a live market for credit risk. If pALPHA trades below NAV, the
market may be pricing future losses or low liquidity. If pPRIME trades at a
premium, the market may be expressing demand for protected exposure. The spread
between NAV and market price becomes information.

\subsection{NAV Premiums and Discounts}

A tranche token can trade above, below, or near NAV.

\begin{center}
\begin{tabular}{>{\bfseries}p{0.22\textwidth}p{0.32\textwidth}p{0.32\textwidth}}
\toprule
Market State & Possible Meaning & User Interpretation \\
\midrule
Premium to NAV & Demand for exposure, expected future income, low perceived risk & Market is paying more than current redemption value \\
Near NAV & Stable expectations, balanced liquidity & Market and accounting value are aligned \\
Discount to NAV & Fear of future losses, low liquidity, urgent exits & Market requires a risk premium \\
\bottomrule
\end{tabular}
\end{center}

The protocol does not force price to NAV because that difference is useful. A
credit market should reveal sentiment, not hide it.

\subsection{Secondary-Market Exit}

Vault redemption is not the only exit path. A user may prefer to sell tranche
tokens to another participant. This can be useful when:

\begin{itemize}
  \item the user wants immediate liquidity,
  \item the vault has a withdrawal queue or liquidity constraint,
  \item the user believes the market price is favorable,
  \item the user wants to rebalance from one tranche to another,
  \item a trader wants to express a view on future credit events.
\end{itemize}

The market layer therefore turns a static credit position into a tradable risk
instrument.

\subsection{Market Makers}

Market makers provide inventory on both sides of a tranche market. They take on
inventory risk in exchange for fees and potential arbitrage opportunities. In a
healthy PRISM market, market makers help connect accounting value to tradable
liquidity while still allowing discounts and premiums to emerge.

\subsection{Market Reaction to Credit Events}

When a credit event occurs, NAV changes first. Traders then decide what the new
information means for future outcomes.

For example:

\begin{enumerate}
  \item A loss event wipes Alpha NAV to zero.
  \item pALPHA may still have liquidity in the market.
  \item Traders sell pALPHA because redemption value has collapsed.
  \item The AMM price moves toward the new risk reality.
  \item Analysts can compare the speed and depth of repricing across tranches.
\end{enumerate}

This makes PRISM more than an accounting system. It becomes a live credit-risk
market.

\section{Why Tranche Tokens Instead of Loan Tokens}

Per-loan tokenization creates fragmented liquidity. Every loan becomes its own
market, with its own risk analysis, buyer base, maturity profile, and trading
depth.

PRISM instead tokenizes pooled risk exposure. A vault can contain one or many
credit assets, but users trade the risk layer rather than the individual loan.
This concentrates liquidity and makes markets easier to understand.

\begin{center}
\begin{tabular}{p{0.3\textwidth}p{0.3\textwidth}p{0.3\textwidth}}
\toprule
Model & Liquidity & User Burden \\
\midrule
Per-loan tokens & Fragmented across many assets & Analyze each loan \\
Tranche tokens & Concentrated by risk layer & Choose risk position \\
\bottomrule
\end{tabular}
\end{center}

\subsection{Diversification and Liquidity}

If a vault contains multiple loans, tranche tokens represent exposure to the
pool's structured risk rather than to one borrower. This allows diversified
credit exposure to trade through a smaller number of liquid instruments.

Per-loan markets are useful when investors want to underwrite individual
borrowers. Tranche markets are useful when investors want exposure to a defined
risk layer across a pool. PRISM chooses the second model because it is more
composable and easier to scale.

\subsection{Due-Diligence Compression}

In a per-loan model, every buyer must understand every asset. In a tranche model,
users can focus on:

\begin{itemize}
  \item the vault strategy,
  \item the originator or sourcing process,
  \item tranche seniority,
  \item historical loss and yield behavior,
  \item market price versus NAV.
\end{itemize}

This does not remove the need for credit analysis. It organizes it around the
pool and capital stack instead of around scattered individual instruments.

\section{Architecture}

PRISM is organized into five conceptual layers.

\subsection{Sourcing Layer}

The Sourcing Layer introduces credit into the system. It connects vaults to
borrower demand, originators, collateral systems, credit attestations, or other
forms of validated yield and credit events.

\subsection{Core Protocol Layer}

The Core Protocol Layer enforces vault accounting, tranche state, NAV updates,
yield distribution, and loss allocation.

\subsection{Risk and Market Layer}

The Risk and Market Layer provides liquidity and price discovery through
tranche-token markets. It allows users to exit or reposition without waiting for
the vault redemption path.

\subsection{Integrity Layer}

The Integrity Layer verifies credit events, preserves accounting correctness,
supports transparent event history, and can integrate optional privacy or
identity modules.

\subsection{Surface Layer}

The Surface Layer exposes vault state to users. It shows tranche NAV, market
price, position value, risk exposure, historical events, and strategy-level
views.

\subsection{Separation of Responsibilities}

PRISM separates concerns deliberately:

\begin{itemize}
  \item The vault should not depend on a specific frontend.
  \item The market should not be able to rewrite vault accounting.
  \item The sourcing layer should not obscure how income and losses are applied.
  \item The integrity layer should verify events without changing economic rules
  arbitrarily.
\end{itemize}

This modularity lets the protocol evolve. New credit originators, collateral
systems, analytics interfaces, or privacy tools can be added around the core
without changing the basic waterfall and cascade model.

\subsection{Conceptual Data Flow}

The high-level flow is:

\begin{enumerate}
  \item Users deposit USDC into a selected tranche.
  \item The vault mints tranche tokens.
  \item Credit exposure generates income or loss.
  \item Income updates NAV through the waterfall.
  \item Loss updates NAV through the cascade.
  \item Tranche tokens can be redeemed at NAV or traded at market price.
\end{enumerate}

This flow is intentionally simple. The sophistication comes from the interaction
between deterministic accounting and free market pricing.

\section{Participants}

\begin{center}
\begin{tabular}{>{\bfseries}p{0.25\textwidth}p{0.65\textwidth}}
\toprule
Participant & Role \\
\midrule
Capital providers & Deposit USDC into Prime, Core, or Alpha and receive tranche tokens. \\
Borrowers or originators & Generate credit demand and return principal, yield, or credit event data. \\
Market makers & Provide liquidity for pPRIME, pCORE, and pALPHA markets. \\
Traders & Buy or sell tranche tokens based on risk expectations. \\
Risk verifiers & Supply or validate credit event information. \\
Interfaces & Present vault state, market prices, and position-level risk to users. \\
\bottomrule
\end{tabular}
\end{center}

\subsection{Capital Providers}

Capital providers are not all the same. A conservative allocator may prefer
Prime, a yield-seeking allocator may prefer Core, and a risk-seeking allocator
may prefer Alpha. PRISM lets each user choose a position in the stack instead of
forcing every participant into the same blended exposure.

\subsection{Borrowers and Originators}

Borrowers and originators are the source of credit demand. They may be
businesses, protocols, real-world asset platforms, payment companies, or
cross-chain collateral providers. PRISM does not require all origination to look
the same. It requires that financial outcomes be translated into clear protocol
events: income, repayment, impairment, default, or recovery.

\subsection{Traders}

Traders are essential because they turn credit information into price. A trader
can buy discounted pCORE, sell risky pALPHA, rotate into pPRIME after a negative
event, or provide liquidity to earn fees. This is how credit sentiment becomes
observable.

\subsection{Risk Verifiers}

Risk verifiers provide or validate information that affects the vault. Their role
is especially important for credit events, collateral status, borrower
performance, and recovery outcomes. The core protocol should treat verified
events as inputs and then apply predefined accounting rules.

\section{Use Cases}

PRISM can support multiple categories of credit markets:

\begin{itemize}
  \item structured pools for real-world credit originators,
  \item invoice and trade-finance vaults,
  \item undercollateralized or partially collateralized borrower facilities,
  \item tokenized credit-risk strategies,
  \item tranche-based yield products,
  \item secondary markets for credit exposure,
  \item risk dashboards and analytics products,
  \item future hedging and protection markets.
\end{itemize}

\subsection{Real-World Asset Vaults}

An originator can group similar credit assets into a PRISM vault. Investors then
choose Prime, Core, or Alpha exposure based on desired risk. Instead of creating
a separate market for every invoice or loan, the originator creates a structured
pool.

\subsection{Trade Finance}

Trade finance often involves short-duration credit, invoices, purchase orders,
and working-capital needs. These cash flows can be pooled and tranched. Prime
capital can fund safer portions, while Alpha capital earns higher yield for
absorbing early losses.

\subsection{Cross-Chain Collateral Credit}

Borrowers may hold valuable assets outside the chain where credit is issued.
PRISM's model can support collateral attestations and cross-chain credit
relationships through the sourcing and integrity layers while keeping tranche
accounting on-chain.

\subsection{Risk Strategies}

Portfolio managers can build strategies around tranche tokens. Examples include:

\begin{itemize}
  \item Prime-only protected credit baskets,
  \item Core yield strategies,
  \item Alpha risk-premium strategies,
  \item dynamic rotation between tranches based on market discount,
  \item market-neutral NAV versus AMM arbitrage strategies.
\end{itemize}

\subsection{Analytics and Research}

Because PRISM emits structured financial events, analysts can build dashboards
showing:

\begin{itemize}
  \item historical NAV,
  \item default frequency,
  \item tranche returns,
  \item market discounts and premiums,
  \item liquidity depth,
  \item realized loss allocation,
  \item originator-level performance.
\end{itemize}

\section{Risk Model}

PRISM does not remove credit risk. It makes credit risk explicit.

Key risks include:

\begin{itemize}
  \item \textbf{Credit risk}: Borrowers may fail to repay.
  \item \textbf{Tranche risk}: Junior tranches can lose value quickly during
  defaults.
  \item \textbf{Liquidity risk}: AMM exit price may diverge from NAV.
  \item \textbf{Oracle and verification risk}: Credit event data must be
  accurate.
  \item \textbf{Market risk}: Traders can reprice assets aggressively after new
  information.
  \item \textbf{Smart contract risk}: Protocol rules must be implemented and
  audited carefully.
\end{itemize}

The goal of PRISM is not to promise safety. The goal is to show where risk sits,
how it moves, and what the market thinks it is worth.

\subsection{Credit Risk}

Credit risk is the risk that the borrower or pool fails to generate expected
cash flows. This is the core risk PRISM exposes. Tranching does not eliminate
credit risk. It allocates credit risk across the capital stack.

\subsection{Subordination Risk}

Subordination defines the order of loss absorption. Alpha is exposed first, Core
second, and Prime last. This makes junior tranches more volatile. It also makes
senior tranches dependent on the existence and size of subordinate capital.

\subsection{Liquidity Risk}

A tranche token may have positive NAV but trade at a discount if market
liquidity is thin or if sellers are urgent. Conversely, a token may trade at a
premium if demand for that exposure is high. Users must understand that market
exit value can differ from redemption value.

\subsection{Information Risk}

Credit markets depend on information. If borrower performance, collateral value,
or default events are reported incorrectly, markets may misprice risk. PRISM's
integrity layer exists to make event verification explicit, but the quality of
inputs remains critical.

\subsection{Concentration Risk}

A vault backed by one borrower behaves differently from a vault backed by many
borrowers. Concentrated vaults can offer clearer exposure but higher idiosyncratic
risk. Diversified vaults reduce single-borrower dependence but may be harder to
analyze at the asset level.

\subsection{Mitigation Principles}

PRISM's risk mitigations are structural rather than promotional:

\begin{itemize}
  \item publish tranche priority before deposit,
  \item make NAV and market price visible,
  \item record income and loss events,
  \item separate vault accounting from market pricing,
  \item use explicit rules for edge cases,
  \item preserve a clear event history,
  \item allow users to choose their risk layer.
\end{itemize}

\section{Governance and Controls}

PRISM's base design does not require discretionary management of user funds.
However, credit systems require controlled parameters and verified events.

\subsection{Parameters}

Vault parameters may include:

\begin{itemize}
  \item tranche target yields,
  \item eligible collateral or originator types,
  \item credit event verification rules,
  \item pause conditions,
  \item market fee settings,
  \item recovery distribution rules.
\end{itemize}

These parameters should be visible before users enter a vault. Changing them
should follow governance or admin controls appropriate to the vault type.

\subsection{Pause and Emergency Controls}

A pause mechanism can protect users during abnormal conditions. A good pause
design should block new risky user actions while still allowing necessary
accounting updates. For example, if a credit event must be recognized, the
protocol should be able to apply it rather than freezing incorrect accounting in
place.

\subsection{Permissioned and Permissionless Vaults}

Different credit markets require different access models. PRISM can support:

\begin{itemize}
  \item permissionless vaults for open participation,
  \item permissioned vaults for regulated or institutional credit,
  \item originator-specific vaults,
  \item curated strategy vaults.
\end{itemize}

The accounting model remains the same. The access layer can differ.

\section{Compliance Posture}

PRISM is designed as infrastructure, not as financial advice. The protocol can
show risk layers, market prices, NAV, and historical outcomes. Users or
interfaces may use that information to make decisions, but the protocol itself
does not decide which tranche a user should choose.

\subsection{Non-Custodial Design}

Users hold tranche tokens in their own wallets. The protocol defines redemption
rules and market mechanics, but token ownership remains portable.

\subsection{Transparency of Risk}

Every tranche should clearly communicate:

\begin{itemize}
  \item its priority in the waterfall,
  \item its position in the loss cascade,
  \item its current NAV,
  \item its current market price,
  \item its exposure to future credit events.
\end{itemize}

This is not only a user-experience principle. It is the foundation of responsible
credit-market design.

\section{Token Philosophy}

PRISM does not require a separate protocol token for its core function. The
native financial instruments of the system are the tranche tokens:

\begin{itemize}
  \item pPRIME,
  \item pCORE,
  \item pALPHA.
\end{itemize}

These represent claims on specific risk layers within a vault. A separate
governance or incentive token is unnecessary for the base credit primitive. If a
future token exists, it should have clear utility, such as insurance backstop
capital or governance over risk parameters, rather than being introduced before
the protocol requires it.

\subsection{Why No Protocol Token Is Needed At The Base Layer}

The protocol's main economic objects are already tokens: pPRIME, pCORE, and
pALPHA. Adding another token too early can confuse the system's purpose. A
credit protocol should first prove that risk can be structured, accounted for,
and priced.

A protocol token may become useful later if it performs a real function. For
example:

\begin{itemize}
  \item absorbing tail risk through an insurance reserve,
  \item governing risk parameters,
  \item backstopping originator failures,
  \item aligning long-term liquidity providers.
\end{itemize}

Until such a function is necessary, tranche tokens are sufficient.

\section{Economic Sustainability}

PRISM can support several protocol-level revenue or sustainability models
without changing the core accounting primitive.

\subsection{Market Fees}

Secondary markets can charge swap fees. These fees compensate liquidity
providers and encourage market depth around tranche tokens.

\subsection{Vault Fees}

Vaults may charge management, origination, or performance fees depending on the
strategy and access model. Any fee should be transparent and applied according
to predefined rules.

\subsection{Risk Infrastructure Fees}

As PRISM expands, verification, analytics, reporting, and risk-management
services may become separate revenue layers. These should support the credit
market rather than distort tranche accounting.

\subsection{Alignment Principle}

The protocol should only extract value when it improves market function:

\begin{itemize}
  \item better liquidity,
  \item better risk verification,
  \item better accounting transparency,
  \item better originator access,
  \item better user decision-making.
\end{itemize}

\section{Ecosystem Extensions}

PRISM is a base primitive. Other systems can build around it.

\subsection{Analytics}

Analytics platforms can index vault events, display historical NAV, track market
discounts, compare originators, and quantify realized loss behavior.

\subsection{Strategy Vaults}

Strategy products can allocate across multiple PRISM tranches, rebalance between
Prime/Core/Alpha, or build diversified baskets of credit exposure.

\subsection{Collateral and Identity Systems}

Borrower identity, collateral verification, cross-chain asset control, and
reputation can live around the sourcing and integrity layers. These extensions
make credit underwriting richer without changing how the core waterfall works.

\subsection{Hedging Markets}

Future protection markets can allow participants to buy or sell default
protection on tranche tokens or vaults. This would create a second layer of
credit-risk pricing around the base tranche markets.

\section{Roadmap}

\subsection{Phase 1: Structured Credit Primitive}

The first phase establishes the core protocol:

\begin{itemize}
  \item tranche vaults,
  \item NAV-based accounting,
  \item yield waterfall,
  \item loss cascade,
  \item tranche tokenization,
  \item secondary markets.
\end{itemize}

The goal of Phase 1 is to prove that the primitive works: users can enter a
risk layer, income can move through the waterfall, losses can move through the
cascade, and markets can price the resulting tokens.

\subsection{Phase 2: Market Expansion}

The second phase expands PRISM into a broader credit-risk marketplace:

\begin{itemize}
  \item multiple vaults and originator types,
  \item richer borrower and collateral integrations,
  \item strategy automation,
  \item analytics and risk scoring,
  \item hedging and protection products.
\end{itemize}

The goal of Phase 2 is to move from a single primitive to a credit-risk
marketplace. This phase deepens liquidity, expands sourcing, and makes PRISM
useful for more than one vault type.

\subsection{Phase 3: Institutional Credit Infrastructure}

The third phase focuses on scale:

\begin{itemize}
  \item permissioned and permissionless vaults,
  \item institutional reporting,
  \item compliance-aware access models,
  \item cross-chain credit markets,
  \item insurance and backstop mechanisms,
  \item deeper liquidity infrastructure.
\end{itemize}

The goal of Phase 3 is to make PRISM a durable base layer for institutional and
open credit markets. This includes better reporting, more sophisticated access
controls, larger pools, and deeper integrations with existing financial
infrastructure.

\section{Illustrative Scenarios}

The following scenarios show how PRISM behaves conceptually.

\subsection{Scenario A: Healthy Credit Pool}

A vault receives regular borrower payments. Prime receives its target income
first, Core receives its target income second, and Alpha receives remaining
income. NAV rises across tranches according to waterfall priority.

In this scenario:

\begin{itemize}
  \item Prime earns steady protected yield,
  \item Core earns higher yield after Prime is satisfied,
  \item Alpha benefits from surplus income,
  \item market prices may trade near or above NAV if confidence is high.
\end{itemize}

\subsection{Scenario B: Mild Stress}

A small impairment occurs but is smaller than Alpha's asset base. Alpha absorbs
the full loss. Core and Prime NAV remain unchanged. pALPHA may trade at a
discount as traders reassess future losses, while pPRIME may trade at a premium
because demand for protected exposure increases.

\subsection{Scenario C: Severe Default}

A large loss wipes Alpha and partially reduces Core. Prime remains protected
because subordinate capital absorbs the loss first. This is the clearest example
of structured credit: every tranche experiences the same credit event, but not
the same economic outcome.

\subsection{Scenario D: Flight to Safety}

After negative credit news, traders sell pALPHA and pCORE while buying or holding
pPRIME. Market prices diverge from NAV. This divergence becomes a live signal of
risk appetite.

\subsection{Scenario E: Recovery}

After a default, some value is recovered. The vault applies recovery according
to predefined rules. NAV improves for eligible tranches, and market prices react
again. This shows that credit events are not only binary defaults; they can also
include restructurings and partial recoveries.

\section{Comparison With Existing Models}

\subsection{Traditional Banks}

Banks pool credit risk internally. Depositors usually receive fixed or low
yield, while the bank keeps much of the spread. Risk sits inside the bank's
balance sheet, and external users have limited visibility.

PRISM externalizes the capital stack. Users choose their risk layer directly and
can observe how income and loss affect that layer.

\subsection{Traditional Structured Finance}

Traditional structured finance uses tranching, but it is often private,
document-driven, and slow to settle. PRISM preserves the useful part of the
model, the capital stack, while making accounting and ownership programmable.

\subsection{Overcollateralized DeFi Lending}

Overcollateralized lending is excellent for liquid collateral and leverage. It
is less suited for borrower types that rely on cash flow, invoices, receivables,
or off-chain collateral. PRISM is designed for credit allocation rather than
automatic liquidation alone.

\subsection{Per-Loan Tokenization}

Per-loan tokenization can be useful for bespoke underwriting, but it fragments
liquidity. PRISM pools credit exposure and tokenizes risk layers, creating fewer
and deeper markets.

\section{Conclusion}

PRISM turns credit from an opaque balance-sheet exposure into a programmable
financial primitive. By splitting pooled credit into transparent tranches,
encoding yield and loss rules on-chain, and allowing each risk layer to trade in
open markets, PRISM makes credit easier to inspect, price, transfer, and build
upon.

The core idea is simple:

\begin{quote}
Credit should not be a black box. It should be structured, transparent, and
market-priced.
\end{quote}

PRISM is the infrastructure for that market.

\end{document}
